\documentclass[11pt]{article}
\usepackage{fontspec}
\usepackage[french]{babel}
\usepackage{amsmath}
\usepackage{amssymb}
\usepackage{amsthm}
\usepackage{mathrsfs}
\usepackage{tikz-cd}
\usepackage{enumitem}

\newcommand{\printedpage}[1]{\marginpar{\footnotesize p.~#1}}
\newcommand{\msp}[1]{\marginpar{\footnotesize [#1]}}

\begin{document}

\printedpage{145}
% head: A. Grothendieck : Dérivateurs, Ch. XII, Édition des Universités de Montpellier II et Paris VII
% foot: Version 30 mars 2003

\msp{page 228}
Il se pose donc la question de la factorisation duale d'une flèche $f : X \longrightarrow Y$ de Cat en
\[
\begin{tikzcd}
X \arrow[r, hook, "W\text{-cofibration}", "i_f"'] & \tilde{Y} \arrow[r, "\in W", "p_f"'] & Y.
\end{tikzcd}
\]

Les résultats de \textsc{Thomason} assurent que ça existe, avec $i$ une immersion ouverte de Dwyer, et $p$ une $W$-fibration - du moins dans le cas $W = W_\infty$. Je vais essayer de le voir directement, en calquant dualement la construction de \textsc{Cartan-Serre}.

Il s'agit donc d'abord de factoriser
\[
X \amalg X \xrightarrow{\ \mathrm{codiag}_X\ } X,
\]
et pour ceci je choisis un intervalle fixe
\[
I \ = \ \Phi_n \ = \ (0 \longrightarrow 1 \longleftarrow 2 \longrightarrow \cdots \longrightarrow 2n-1 \longleftarrow 2n)
\]
avec
\[
n \geq 2
\]
de sorte que l'inclusion
\[
\varepsilon = \text{catégorie discrète } \{0,1\} \xrightarrow{\partial} I, \qquad \partial(0) = 0, \quad \partial(1) = 2n,
\]
soit une immersion ouverte de Dwyer. Cette propriété est stable par changement de base Cat-fibrant, en particulier en prenant un produit
\[
\underbrace{X \times \varepsilon}_{\simeq X \coprod X} \xrightarrow{i_X} X \times I,
\]
% continuation of a footnote begun on the preceding page of the edition (no footnote mark on this page)
\begingroup\renewcommand{\thefootnote}{}\footnotetext{%
\[
\begin{tikzcd}[ampersand replacement=\&]
U \arrow[r, hook, "i"] \arrow[d, "p"'] \& X \arrow[d, "q"] \\
U' \arrow[r, hook, "i'"] \& X'
\end{tikzcd}
\]
avec $i$ immersion ouverte de Dwyer, alors
\[
\begin{aligned}
i \in W \quad &\Longleftrightarrow \quad i' \in W \\
p \in W \quad &\Longleftrightarrow \quad q \in W,
\end{aligned}
\]
j'ai donc besoin de W(7 bis) \textit{et} même de W(7) pour la deuxième implication.
\begin{enumerate}
\item[c)] Les immersions ouvertes de Dwyer \texttt{[sont]} stables par cochangement de base et par composition, contiennent les isomorphismes.
\end{enumerate}}\endgroup
\printedpage{146}
% head: A. Grothendieck : Dérivateurs, Ch. XII, Édition des Universités de Montpellier II et Paris VII
% foot: Version 30 mars 2003
\msp{page 229}
d'autre part
\[
X \times I \xrightarrow{\mathrm{pr}_1} X
\]
est $W_\omega$-parfait, d'où la factoriastion cherchée % [sic] factoriastion -> factorisation, misprint
\[
X \amalg X \xrightarrow{i_X} X \times I \ \xrightarrow{\ p_X = \mathrm{pr}_1^X\ } X
\]
- il n'y a pas plus simple! Comment l'utiliser pour construire la factorisation d'une flèche $f : X \longrightarrow Y$ quelconque, en dualisant les constructions relatives aux espaces de chemins? On définit $Z(f)$ par le carré cocartésien dans Cat
\[
\begin{tikzcd}
X \arrow[r, dotted, "\mathrm{inj}_1^X"'] & X \times I \arrow[r, "f'"] & Z(f) \arrow[d, dotted, bend left, "p_f"] \\
& X \arrow[u, hook, "\mathrm{inj}_0^X"'] \arrow[r, "f"] & Y, \arrow[u, hook, "W"']
\end{tikzcd}
\]
on a
\[
Z(f) \xrightarrow{\alpha} Y \qquad \text{induit par } \mathrm{pr}_1 : X \times I \longrightarrow X,
\]
et la factorisation cherchée de $f$ est
\[
X \ \xrightarrow{\ i_f = f' \circ \mathrm{inj}_1^X\ } Z(f) \xrightarrow{p_f} Y.
\]

Je dis que
\[
\left\{
\begin{aligned}
& i_f \text{ est immersion ouverte de Dwyer, donc une } W\text{-cofibration.} \\
& p_f \in W.
\end{aligned}
\right.
\]

\msp{page 230}
On voit sur le diagramme plus complet
\[
\begin{tikzcd}
& X \times I \arrow[r, "f'"] \arrow[dd, dotted, bend left=20, "\mathrm{pr}_1"] \arrow[dr, phantom, "\substack{\text{cocart.}\\ \text{et cart.}}"] & Z(f) \arrow[dd, dotted, bend left=60, "p_f"] \\
X \arrow[ur, dotted, "\mathrm{inj}_1^X"] \arrow[r, hook, "i_1^X"] & X \amalg X \arrow[u, hook, "i_X"'] \arrow[r, "f \coprod \mathrm{id}_X"] \arrow[dr, phantom, "\substack{\text{cocart.}\\ \text{et cart.}}"] & Y \amalg X \arrow[u, hook, "j"'] \\
& X \arrow[u, hook, "i_0^X"'] \arrow[r, "f"] & Y, \arrow[u, hook]
\end{tikzcd}
\]
\printedpage{147}
% head: A. Grothendieck : Dérivateurs, Ch. XII, Édition des Universités de Montpellier II et Paris VII
% foot: Version 30 mars 2003
\[
i_f = f' \circ \mathrm{inj}_1^X = f' i_X i_1^X = j \ \ \underbrace{(f \amalg \mathrm{id}_X) i_1^X}_{\substack{\text{inclusion canonique}\\ \text{de } X \text{ dans}\\ Y \coprod X}} \ \ ,
\]
c'est \texttt{[une]} immersion ouverte de Dwyer, comme composé de $Y \longrightarrow Y \amalg X$ et $j : Y \amalg X \longrightarrow Z(f)$ qui le sont ($j$ obtenu par cochangement de base à partir de $i_X$, qui est immersion ouverte de Dwyer).

D'autre part, $X \xrightarrow{\mathrm{inj}_0^X} X \times I$ est \texttt{[une]} immersion ouverte de Dwyer et dans $W$, donc coüniversellement dans $W$, donc $Y \longrightarrow Z(f)$ est elle aussi dans $W$, donc aussi la rétraction $p_f$.

D'autre part, la proposition duale de
\[
\underline{\mathrm{Fib}}_W S \cap W_S = W^{\mathrm{u}}_{S\,\mathrm{univ}}
\]
est valable pour les $W$-cofibrations :

\msp{page 231}
\textbf{Proposition 16} (\footnote[125]{N'a rien à voir avec Cat - une chose générale quand il y a assez de $W$-cofibrations.}). \textit{Sur la sous-catégorie pleine $\underline{\mathrm{Cofib}}_W(S)$ de $S\backslash\mathrm{Cat}$, le localiseur ${}_SW$ est égal au $({}_SW)^{\mathrm{univ}}$ ($W$-équivalences couniverselles), i.e. on a ceci : Soient}
\[
S \longrightarrow X, \qquad S \longrightarrow Y
\]
\textit{des $W$-cofibrations, et soit $f : X \longrightarrow Y$ un morphisme de $S\backslash\mathrm{Cat}$ tel que $f \in W$. Alors pour tout cochangement de base $S \longrightarrow S'$, la flèche correspondante}
\[
f' : X' = X \amalg_S S' \longrightarrow Y' = Y \amalg_S S'
\]
\textit{est dans $W$.}

\textsc{Démonstration}. On factorise $S \longrightarrow S'$ en
\[
S \xrightarrow{i} S'_0 \xrightarrow{j} S',
\]
avec $i$ immersion ouverte de Dwyer (donc $W$-cofibration), et $j \in W$. On a donc le diagramme à carrés cocartésiens
\[
\begin{tikzcd}
S \arrow[r] \arrow[d, hook, "i"'] & X \arrow[r, "f"] \arrow[d, hook] & Y \arrow[d, hook] \\
S'_0 \arrow[r] \arrow[d, "\in W"'] & X'_0 \arrow[r, "f'_0", "W"'] \arrow[d, "W"] & Y'_0 \arrow[d, "W"] \\
S' \arrow[r] & X' \arrow[r, "f'"] & Y'.
\end{tikzcd}
\]
Comme $i$ est une $W$-cofibration et $f \in W$, on a $f'_0 \in W$. D'autre part, comme $S \longrightarrow X$, donc $S'_0 \longrightarrow X'_0$ est
\printedpage{148}
% head: A. Grothendieck : Dérivateurs, Ch. XII, Édition des Universités de Montpellier II et Paris VII
% foot: Version 30 mars 2003
\msp{page 232}
une $W$-cofibration, $X'_0 \longrightarrow X'$ est dans $W$, et pour la même raison, $Y'_0 \longrightarrow Y$ l'est aussi. Donc $f'$ l'est aussi. % [sic] Y'_0 \longrightarrow Y -> Y'_0 \longrightarrow Y', the diagram has Y'_0 \to Y'

\textbf{NB} Ici, c'est un argument formel, valable (sous forme duale) dans tout $(\mathcal{M}, W)$ où on a un énoncé de factorisation pour une flèche quelconque $f : X \longrightarrow Y$ en $X \xrightarrow{i_f} \bar{X} \xrightarrow{\alpha} Y$, avec $i_f \in W$ et $\alpha$ une $W$-fibration. Alors, si $X$, $Y$ sont $W$-fibrés sur $S$, toute $S$-flèche de $X$ dans $Y$ qui est dans $W$, est dans $W^{\mathrm{u}}_S$. Et dualement, si toute flèche $X \longrightarrow Y$ se factorise en $X \longrightarrow \tilde{Y} \longrightarrow Y$, avec $X \longrightarrow \tilde{Y}$ une $W$-cofibration et $\tilde{Y} \longrightarrow Y$ dans $W$.

Il me vient ici une (petite) inquiétude : ai-je au moins prouvé que si
\[
S \longrightarrow X
\]
est une $W$-cofibration, alors pour toute flèche $S' \longrightarrow S$ dans Cat, le carré cocartésien dans Cat % [sic] S' \longrightarrow S -> S \longrightarrow S', the square below has S \to S'
\[
\begin{tikzcd}
S \arrow[r] \arrow[d] & X \arrow[d] \\
S' \arrow[r] & X'
\end{tikzcd}
\]
est aussi $W$-cocartésien (\footnote[126]{dans le sens ancien, spécial à Cat (déf. 3, p. 131). Pour le comparer avec la théorie générale, cf. prop. 20, p. 262.})? Autrement,
\msp{page 233}
à quoi ce carré-là serait-il bon?! Pour le prouver, factorisons encore $S \longrightarrow S'$ en
\[
\underset{\substack{\text{immersion}\\ \text{ouverte de}\\ \text{Dwyer}}}{S \hookrightarrow} S'_0 \underset{\in W}{\longrightarrow} S',
\]
d'où le diagramme de 2 carrés cocartésiens
\[
\begin{tikzcd}
S \arrow[r, "\mathrm{Cof}_W"] \arrow[d, hook, "D_W"'] & X \arrow[d, hook, "D_W"] \\
S'_0 \arrow[r, "\mathrm{Cof}_W"] \arrow[d, "\in W"'] & X'_0 \arrow[d, "\in W"] \\
S' \arrow[r, "\mathrm{Cof}_W"] & X',
\end{tikzcd}
\]
où $X'_0 \longrightarrow X'$ est $\in W$ d'après l'hypothèse sur $X$ sous $S$. Il faut donc vérifier :
\printedpage{149}
% head: A. Grothendieck : Dérivateurs, Ch. XII, Édition des Universités de Montpellier II et Paris VII
% foot: Version 30 mars 2003
\begin{enumerate}
\item[a)] Un carré cocartésien construit sur une flèche qui est \texttt{[une]} immersion de Dwyer, est $W$-cocartésien.
\item[b)] Un carré commutatif dans lequel deux côtés opposés sont dans $W$, est $W$-cocartésien.
\item[c)] Le composé de deux carrés $W$-cocartésiens est $W$-cocartésien.
\end{enumerate}

Admettons provisoirement b) et c), et prouvons a), pour un carré cocartésien
\[
\begin{tikzcd}
U \arrow[r, hook, "i"] \arrow[d, "p"'] & X \arrow[d, "q"] \\
U' \arrow[r, hook, "i'"] & X',
\end{tikzcd}
\]
où $i$ est immersion ouverte de Dwyer. On procède encore comme tantôt, en factorisant $p$ en $U \overset{j}{\hookrightarrow} U'_0 \xrightarrow{p'} U'$, où $j$ est \texttt{[une]} immersion ouverte de Dwyer, $p' \in W$. On a donc un diagramme de deux carrés cocartésiens accolés
\msp{page 234}
\[
\begin{tikzcd}
U \arrow[r, hook, "i \in D_W"] \arrow[d, hook, "D_W"'] & X \arrow[d, hook] \\
U'_0 \arrow[r, hook, "D_W"] \arrow[d, "p' \in W"'] & X'_0 \arrow[d, "q' \in W"] \\
U' \arrow[r, hook] & X',
\end{tikzcd}
\]
où le premier carré est formé d'immersions ouverts. On sait alors qu'il est $W$-cocartésien. Dans le deuxième, on a $p' \in W$, donc $q' \in W$ puisque $i$ est une immersion de Dwyer, donc une $W$-cofibration. Donc il est également $W$-cocartésien. On conclut alors grâce à c) ci-dessus. % [sic] immersions ouverts -> immersions ouvertes, agreement

\noindent Prouvons b), pour
\begin{tikzcd}[baseline=(\tikzcdmatrixname-1-1.base)]
X_0 \arrow[r, "i_1"] \arrow[d, "i_2 \in W"'] & X_1, \arrow[d, "f_1 \in W"] \\
X_2 \arrow[r, "f_2"] & Y
\end{tikzcd}
en l'insérant dans
\begin{tikzcd}[baseline=(\tikzcdmatrixname-1-1.base)]
X_0 \arrow[r, "i_1"] \arrow[d, "i_2 \in W"'] & X_1, \arrow[d, "j_1"] \arrow[ddr, "f_1"] & \\
X_2 \arrow[r, "j_2"] \arrow[drr, "f_2"'] & \mathcal{X} \arrow[dr, "\varphi"] & \\
& & Y
\end{tikzcd}
où $\mathcal{X} = \displaystyle\int X_0$
\begin{tikzcd}[baseline=(\tikzcdmatrixname-2-1.base), row sep=small]
& X_1 \\
{} \arrow[ur] \arrow[dr] & \\
& X_2
\end{tikzcd}
.

Comme $i_2 \in W$, on sait qu'il en résulte $j_1 \in W$, et comme $f_1 \in W$, il en résulte que $\varphi \in W$, q.e.d.

\end{document}
